\hypertarget{hem-survives-scrutiny-as-the-starting-point}{%
\section{HEM Survives Scrutiny as the Starting
Point}\label{hem-survives-scrutiny-as-the-starting-point}}

\textbf{Bottom line: the Homogeneous Equilibrium Model is a defensible
--- indeed, the literature-optimal --- first extension of the existing
single-phase solver, but only with three explicit caveats attached.} The
quantitative case is strong: fitted CO₂ phase-change relaxation times
are \textbf{0.04--2.82 ms} against a wellbore blowdown time scale of
minutes (equilibrium Damköhler number ≳10⁵), the only controlled
vertical two-phase CO₂ dataset found near-zero slip at 6.5 MPa, and
SINTEF's own six-equation two-fluid model ``did not clearly outperform
the HEM'' against pipe-depressurization data. The caveats: (i) HEM
systematically underpredicts choked discharge through short restrictions
by 0--26\% --- which is why the choke boundary, not the wellbore
interior, should carry the non-equilibrium correction; (ii) the HEM
mixture sound speed is discontinuous at the saturation boundary, the
principal numerical hazard for the HLLC solver; and (iii) the no-slip
assumption fails precisely in the late-blowdown, low-rate tail, where
condensate fallback in a 17 cm bore at an area ratio of
\textasciitilde200:1 to the 12.7 mm choke is governed by drift
velocities of order 0.1--1 m/s that exceed the mixture velocity. The
staged path HEM → HRM → drift-flux is mathematically protected by the
subcharacteristic condition (each added equilibrium assumption lowers
the sound speed, so the equilibrium limit is a well-behaved limit of the
relaxation system), and the whole architecture --- MUSCL-Hancock/HLLC,
Span--Wagner tabulated closure, conjugate radial heat transfer,
quasi-steady choke boundary --- survives the extension essentially
intact. The deepest finding of this review is negative and valuable:
\textbf{no vertical CO₂ blowdown dataset, no published CO₂-well DTS
phase-transition trace, no measured condensation-side relaxation time,
and no transport closure validated above \textasciitilde10 mm diameter
exist} --- four publishable gaps the planned PERTT experiment sits
directly on.

\begin{center}\rule{0.5\linewidth}{0.5pt}\end{center}

\hypertarget{the-literature-splits-into-pipeline-decompression-choke-discharge-and-a-thin-wellbore-corpus}{%
\subsection{1. The literature splits into pipeline decompression, choke
discharge, and a thin wellbore
corpus}\label{the-literature-splits-into-pipeline-decompression-choke-discharge-and-a-thin-wellbore-corpus}}

The transient CO₂ depressurization literature has three centers of
gravity --- SINTEF/NTNU (pipe depressurization, relaxation-model theory,
choked orifice flow), UCL/Mahgerefteh (pipeline rupture, HRM, solids)
--- and a smaller wellbore corpus (LBNL/T2Well, SINTEF well models,
OLGA). The per-study record below carries forward the research notes'
verification flags: references marked \textbf{{[}unverified{]}} were not
confirmed against a primary source and must be checked before citation.

\hypertarget{pipeline-and-shock-tube-decompression}{%
\subsubsection{1.1 Pipeline and shock-tube
decompression}\label{pipeline-and-shock-tube-decompression}}

\textbf{\citet{munkejord2020} --- the anchor dataset.} ECCSEL rig: 61.67 m
SS316 horizontal tube, 40.8 mm ID, honed bore, 16 Kulite transducers
logged at 100 kHz concentrated near the rupture end, 23 Type-E
thermocouples (±0.22 °C) including top/side/bottom triplets at four
stations, full-bore rupture-disc opening; pressure uncertainty 60 kPa
(95\%). Tests span gas-phase (Test 3: \textbf{4.04 MPa, 10.2 °C} --- the
closest published analog to the 2,000-ft well's initial state) to dense
phase (12.54 MPa/21.1 °C; 10.40 MPa/40 °C), plus an N₂ single-phase
reference. Model: HEM with Span--Wagner EOS. Conclusions: single-phase
wave speeds match Span--Wagner ``very well''; measured two-phase plateau
pressures fall \textbf{below} the equilibrium prediction (delayed
flashing); the equilibrium wave-speed kink (e.g., \textasciitilde370 m/s
→ \textless40 m/s at the phase boundary in Test 8) is measured as a
smooth transition lying between equilibrium and frozen curves; HEM
predicts dry-out \textasciitilde0.5 s early; in the gas-phase Test 3,
condensation droplets appear near \textbf{2.8 MPa, below the equilibrium
dew pressure} --- delayed condensation, the gas-side analog of delayed
flashing; the N₂ closed-end comparison exposed a
heat-transfer-correlation failure at low velocity (−150 °C predicted vs
−78 °C measured), not a phase-equilibrium failure. Raw data are open \citep{sintef2020data} (\href{https://www.pvv.ntnu.no/~stm/dring/depress-co2-exp.pdf}{preprint};
\href{https://sintef.brage.unit.no/sintef-xmlui/handle/11250/2770324}{SINTEF
Brage}).

\textbf{\citet{munkejord2021}} --- same rig, CO₂--N₂ and CO₂--He;
isoenergetic--isochoric flash via an in-house VLE algorithm interfaced
to TREND; decompression-wave speed ``abruptly decreases from about 350
m/s to 40 m/s at 6 MPa'' under full equilibrium; the production numerics
are the dissipative FORCE flux + MUSCL + SSP-RK2 at CFL 0.85, 1200--4800
cells; open data \citep{sintef2021data} (\href{https://www.pvv.ntnu.no/~stm/dring/depress-mix-exp.pdf}{preprint}).

\textbf{\citet{log2024hrm}} --- HEM vs HRM* (constant,
per-case-fitted relaxation time) across initial temperatures 10--40 °C:
the fitted relaxation time \textbf{decreases with initial temperature
and spans \textasciitilde60×}, 0.04 ms (40 °C) / 1.50 ms (21 °C) / 2.82
ms (4.6 °C)
(\href{https://sintef.brage.unit.no/sintef-xmlui/handle/11250/3103647}{Brage
record}).

\textbf{\citet{log2024hfm} (HFM)}
--- a four-equation homogeneous model plus
bubble-number/interfacial-area transport, CNT homogeneous nucleation
plus fitted heterogeneous rates (``not predictive''), kinetic-theory
evaporation; one global parameter set achieves MAPE \textbf{6.0\% vs
7.6\%} for per-case-tuned HRM*; numerics are Godunov splitting + HLLC +
implicit Euler sources. Metastability findings: near-critical dome
entries (1.2 K below T\_c) show essentially zero undershoot; five colder
tests flashed where CNT predicts \textless10⁻¹⁰ bubbles m⁻³s⁻¹,
implicating un-modeled heterogeneous nucleation; crack-tip pressures
reach 20\% below saturation
(\href{https://www.pvv.ntnu.no/~stm/dring/hfm_preprint.pdf}{HFM
preprint}).

\textbf{\citet{munkejord2015}} --- HEM vs a six-equation two-fluid model
(both thermodynamic-equilibrium; slip via interfacial + wall friction),
Peng--Robinson + TRAPP transport, validated on National Grid Test 31
\citep{cosham2012}, Botros CO₂--CH₄, Drescher CO₂--N₂. Dry-out timing
errors of opposite sign (HEM \textasciitilde2 s early, TFM
\textasciitilde4 s late on the 90\% CO₂ case); the authors conclude the
TFM \textbf{``did not clearly outperform the HEM''} and prefer HEM for
simplicity, noting better TFM closures would require local experimental
data that do not exist
(\href{https://www.sintef.no/globalassets/project/co2-dynamics/publications/depress-mix-preprint.pdf}{preprint}).

\textbf{\citet{brown2013}} --- HRM for full-bore rupture of
dense-phase CO₂ pipelines with an empirically derived relaxation-time
equation; verified on hypothetical shock tubes, validated against
CO2PipeHaz full-bore ruptures. The load-bearing conclusion:
\textbf{delayed phase transition has negligible effect on the
decompression-wave rate, but ignoring it underestimates the transient
discharge rate}
(\href{https://www.scipedia.com/public/Brown_et_al_2013a}{Scipedia};
\href{https://hal-ineris.archives-ouvertes.fr/ineris-00963485}{HAL-INERIS}).
The exact relaxation-time correlation and EOS were not retrievable from
open pages \textbf{{[}Downar-Zapolski-form refit + Peng--Robinson
recalled from training, unverified{]}}.

\textbf{\citet{martynov2018}} --- 36.7 m, 50 mm pipe, ≤45 bar, 4/6 mm
orifices: pressure and temperature stabilize near the triple point;
video shows liquid/vapour \textbf{stratification then rapid solid-CO₂
formation}; an HEM-based model with vapor--liquid--solid properties
predicts P(t), T(t) and solid mass, self-described as ``partially
validated''
(\href{https://eprints.whiterose.ac.uk/id/eprint/130497}{White Rose}).
The observed stratification before solidification is direct evidence
that homogeneous assumptions degrade near the triple point.

\textbf{Botros et al.~shock tubes} (decompression-wave speed vs
pressure, the fracture-control community's standard): \citet{botros2010}; \textbf{\citet{botros2016} (the key pure-CO₂ paper)}; plus
binary-impurity companions in JPVT 139(2) (DOI 10.1115/1.4034016,
10.1115/1.4034466) and \emph{Can. J. Chem. Eng.} 95(3):440--448 (DOI
10.1002/cjce.22689) --- all DOIs verified from the reference list of \citet{munkejord2021}
(\href{https://www.pvv.ntnu.no/~stm/dring/depress-mix-exp.pdf}{mirror}).
Rig dimensions (\textasciitilde42 m, \textasciitilde38 mm ID) are
\textbf{{[}unverified{]}}.

\textbf{COOLTRANS (National Grid/Spadeadam)} --- \citet{cosham2012}: 144 m × 150 mm shock tube; pressure ``fell to saturation quickly, with
little variation along the 144 m length.'' The vertical venting rig
(dense CO₂ at 150 barg through 12/25/50 mm vertical vents, Coriolis mass
flow) found predictions matching the smaller vents while the \textbf{50
mm vent out-flowed homogeneous-equilibrium predictions --- delayed
nucleation}; small amounts of solid CO₂ were observed in vents and
craters
(\href{https://www.icheme.org/media/9161/paper20-hazards-23.pdf}{Allason
et al., Hazards XXIII}).

\textbf{Guo et al.~(DUT/UCL, CO2QUEST)} --- 258 m × 233 mm
industrial-scale pipeline discharging through a \textbf{vertical 15 mm
orifice}; Test 1 is \textbf{gaseous CO₂ at 3.7 MPa, 37.5 °C} --- the
single large-scale test closest to the well's initial state, including
the vertical exit geometry; 10 pressure sensors at 100 kHz, 24
thermocouples resolving top/bottom stratification
(\href{https://discovery-pp.ucl.ac.uk/id/eprint/10051724/1/Manuscript_01.pdf}{UCL
manuscript}). Companion journal papers (Applied Energy 178:189--197;
Applied Energy 183:1279--1291; Energy 119:53--66) are cited from that
manuscript's reference list without DOIs.

\textbf{Vessel blowdowns}: \citet{gebbeken1996} supercritical-CO₂ vessel blowdown with void fraction and stratification is the classical reference but its full citation \textbf{could not be verified this session {[}unverified --- training recollection{]}}. A 2025 open dataset exists for a 2 m³ sphere to
130 bar and −50 °C (INERIS;
\href{https://www.ncbi.nlm.nih.gov/pmc/articles/PMC12819109/}{PMC12819109},
title and PMC ID only --- retrieve before use).

\hypertarget{chokeorifice-critical-flow}{%
\subsubsection{1.2 Choke/orifice critical
flow}\label{chokeorifice-critical-flow}}

\textbf{\citet{hammer2022} --- the direct precedent for the planned architecture.} Six
pipe blowdowns from \textasciitilde12 MPa/25 °C through orifices and
nozzles of 4.5, 9.0 and \textbf{12.7 mm} (the exact target choke size),
coupled to a 1-D transient HEM pipe solver (4800 cells, MUSCL+RK,
Friedel friction). Findings: \textbf{``the choke point is at a
non-equilibrium state''}; HEM mass-flux error −26\% to 0\% (always
underpredicting), Henry--Fauske −22\% to +5\%, delayed-HEM (flashing
deferred to the CNT superheat limit) −21\% to +10\%; against
Banasiak--Hafner nozzle data, MARE 7.8\% (HEM), 2.5\% (D-HEM), 5.4\%
(HF), with HEM degrading to 17\% on the low-pressure subset. Orifice
contraction coefficients computed from a Bragg-type force balance:
\textbf{C\_c = 0.75 (12.7 mm), 0.74 (9.0 and 4.5 mm)}, from the
incompressible sharp-edge limit π/(π+2) ≈ 0.611; the authors state the
orifice value remains uncertain. HEM's choke search ``often'' locks at
the saturation crossing because the sound speed is discontinuous there
(\href{https://sintef.brage.unit.no/sintef-xmlui/handle/11250/3023978}{SINTEF
record};
\href{https://www.pvv.ntnu.no/~stm/dring/outflow_preprint.pdf}{preprint}).

\textbf{Boccardi et al.~(ENEA)} --- 107 steam-water PSV tests (6 mm
orifice, 0.5--1.75 MPa, inlet quality 0.5--10\%): HEM never
overestimated, underpredicting 2--40\%; HNE-DS variance 6.9\%;
Henry--Fauske mean −0.25\% but overestimating up to 27.4\%
(\href{https://www.aidic.it/CISAP4/webpapers/33Boccardi.pdf}{CISAP4
paper}). Fauske's relaxation-length framework (non-equilibrium vanishes
beyond \textasciitilde0.1 m flow length) is documented in a LANL review
(\href{https://public.lanl.gov/davidp/MSThesisHTML/2.html}{LANL thesis
review}). \citet{henry1971} and the N = x\_eq/0.14 form are \textbf{{[}unverified --- check
original{]}}; \citet{morris1990} reports
±6--7\% reproduction of choked gas-liquid flow by empirical C\_d
correlations
(\href{https://researchportal.hw.ac.uk/en/publications/discharge-coefficients-for-choked-gas-liquid-flow-through-nozzles/}{portal}).

\hypertarget{wellbore-flow}{%
\subsubsection{1.3 Wellbore flow}\label{wellbore-flow}}

\textbf{T2Well/ECO2N \citep{pan2011}} --- 1-D
transient non-isothermal \textbf{drift-flux} CO₂--water--NaCl wellbore
model coupled to TOUGH2/ECO2N: fully implicit mass/energy (Newton),
semi-explicit mixture momentum; \citet{shi2005} slip closure with
profile parameter C₀ capped 1.0--1.5 and a Kutateladze-number drift
velocity re-tuned to Ku = 142 (vs Richter's 75); Ramey or gridded radial
heat exchange; explicitly \textbf{cannot represent coexisting liquid +
gaseous CO₂ with a third phase}; drift velocities 0.28--0.72 m/s in its
verification case; validation limited to analytics and leakage scenarios
--- no field two-phase data
(\href{https://escholarship.org/content/qt0k35m2vg/qt0k35m2vg.pdf}{manual
PDF}; \href{https://www.osti.gov/biblio/1007233}{OSTI 1007233}). \citet{shi2005} itself was obtained only second-hand
\textbf{{[}forms must be taken from the original{]}}.

\textbf{SINTEF two-fluid well model \citep{linga2016}} --- two-fluid vertical-well model with
conjugate conduction through tubing/annulus/casing/cement/rock; phase
mass transfer via \textbf{relaxation source terms that satisfy the
second law and are continuous in the single-phase limit} (engineered
specifically for a robust single↔two-phase transition);
flow-regime-dependent friction and heat transfer; Span--Wagner EOS;
applied to shut-in and blowout
(\href{https://sintef.brage.unit.no/sintef-xmlui/handle/11250/2620521}{Brage
record}; author attribution to Linga \& Lund \textbf{unverified}).

\textbf{UCL well model (Samuel \& Mahgerefteh)} --- 1-D HEM,
conservative hyperbolic form, Godunov FV, Peng--Robinson EOS, Chen
friction, Dittus--Boelter in-well heat transfer, wellhead ghost cells
with isentropic/Thompson characteristic relations; \textbf{validated
against Ketzin single-phase wellbore temperatures} (0.41 kg/s; Möller et
al.~2014 data) and benchmarked against OLGA, where PR gave
systematically lower temperatures --- a cubic-EOS bias argument for
Span--Wagner
(\href{https://discovery-pp.ucl.ac.uk/10080131/3/Samuel\%20Ramp-up\%20injection\%20paper.pdf}{UCL
PDF}). Their Goldeneye ramp-up cases show slower ramps giving
\emph{colder} wellhead minima (252/238/226 K for 5 min/30 min/2 h
ramps).

\textbf{\citet{lindeberg2011} (SINTEF, TCCS-6)} --- quasi-steady marching
blowout model of the Sleipner well: rates 181--313 kg/s, outlet
temperatures −75 to −97 °C, hydrate zones 200--550 m below wellhead with
plug-melt \textbf{pulsed release}; dry ice in the well judged unlikely;
explicitly called for an instrumented abandoned-well blowout experiment
(\href{https://www.sintef.no/globalassets/project/tccs6/presentations/e6/03_co2blowout.pdf}{TCCS-6
PDF}). His adiabatic-well conclusion applies to sustained high-rate
blowouts and \textbf{does not transfer} to a slow, finite-inventory,
choked blowdown where conjugate wall heat input is first-order.

\textbf{\citet{hammer2021} (preprint; journal of record unconfirmed)} --- the only controlled vertical
two-phase CO₂ dataset: IFE FALCON loop, \textbf{44 mm vertical pipe at
6.5 MPa}, superficial velocities 0.2--4.0 m/s (gas) / 0.15--3.0 m/s
(liquid), holdup by X-ray and gamma densitometry. Result: \textbf{``the
experiments show nearly no-slip flow, except at low inlet liquid
holdup,'' and ``the simplest model, the homogeneous no-slip approach,
performed best''} --- beating both \citet{shi2005} and the T2Well adaptation, because gas and liquid properties converge near the critical
point
(\href{https://www.pvv.ntnu.no/~stm/dring/driftflux-well.pdf}{preprint}).

\textbf{Field/industry context}: \citet{wildwell2024} (Havtil note 2024-380) documents dry ice ejected from wellbores, frozen annuli at
\textasciitilde100 m casing shoes, and the judgment that with CO₂'s
critical point inside typical operating envelopes ``a phase change is
almost guaranteed somewhere in the wellbore if control is lost''; the
new OLGA CO₂ solver tightens p--T coupling for stability but is
proprietary
(\href{https://www.havtil.no/contentassets/9b77c8de198b49afb887f9688b45f4be/havtil_co2-blowoutconsiderations_techncialnote_2024-380.pdf}{Havtil
PDF}).

\textbf{Negative results carried forward explicitly}: no measured
in-well dataset from any CO₂ well blowdown exists in the open
literature; no DTS record of a two-phase CO₂ column tracking T\_sat(p)
was found (Ketzin deliberately avoided two-phase operation with N₂
flushes, keeping \textasciitilde3 °C margin to the condensation line ---
\href{https://iahs.info/uploads/dms/15943.50-286-290-355-40-Wiese.pdf}{Wiese
et al.~2012};
\href{https://gfzpublic.gfz-potsdam.de/pubman/item/item_242409_1/component/file_242408/15975.pdf}{Würdemann
et al.~2010}); no CO₂-specific CCFL/flooding study for vertical
wellbores was found; Ruan et al.~(IJGGC 2013) and Lu \& Connell wellbore
papers \textbf{could not be located --- do not cite}.

\hypertarget{model-comparison-the-relaxation-hierarchy-orders-the-choices-and-the-data-refuse-to-reward-complexity}{%
\subsection{2. Model comparison: the relaxation hierarchy orders the
choices and the data refuse to reward
complexity}\label{model-comparison-the-relaxation-hierarchy-orders-the-choices-and-the-data-refuse-to-reward-complexity}}

The mathematically clean frame is the relaxation hierarchy of \citet{lund2012} (\href{https://www.sintef.no/globalassets/project/co2-dynamics/publications/lund_two-phase_relaxation_hierarchy.pdf}{SINTEF PDF}; journal metadata \textbf{{[}from training, unverified{]}}): a homogeneous
(single-velocity) two-phase model with three relaxation processes ---
volume transfer driven by Δp, heat by ΔT, mass by Δμ --- generates eight
models; imposing each equilibrium \textbf{lowers} the mixture sound
speed through sum-of-squares increments (1/ã²\_X = 1/ã²\_Y + Z ≥ 0), so
every equilibrium limit satisfies the \textbf{subcharacteristic
condition} for any thermodynamically stable EOS. HEM (p--T--μ
equilibrium) is the bottom of the hierarchy with the lowest sound speed
and is the \textbf{only} member whose sound speed is discontinuous at
the single-phase boundary. HRM is the hierarchy's pT model with finite
μ-relaxation --- strictly dissipative relative to HEM and smooth at the
dome boundary. Velocity non-equilibrium (drift-flux, two-fluid) sits
outside the proven ordering; Lund defers it explicitly.

Against this frame, the empirical record is one-directional. HEM with a
reference EOS captures decompression-wave speeds well and bulk two-phase
behavior ``fairly'' but overpredicts the pressure plateau for dense
initial states and underpredicts short-restriction discharge; HRM* fixes
the plateau at the cost of a case-dependent relaxation time (0.04--2.82
ms, \textasciitilde60× spread); the nucleation-based HFM achieves
HRM*-level accuracy (MAPE 6.0\% vs 7.6\%) with one global parameter set
but fitted heterogeneous rates; the six-equation two-fluid model did
\textbf{not} beat HEM and its own authors recommend HEM
(\citealp{munkejord2015}, \href{https://www.sintef.no/globalassets/project/co2-dynamics/publications/depress-mix-preprint.pdf}{preprint}; \citealp{log2024hfm}, \href{https://www.pvv.ntnu.no/~stm/dring/hfm_preprint.pdf}{HFM preprint}). Drift-flux earns its keep only where mixture velocity falls
to the drift-velocity scale --- O(0.1--1 m/s) per T2Well's Kutateladze
closure --- and the one vertical CO₂ dataset shows no-slip winning
outright at 6.5 MPa
(\citealp{hammer2021}, \href{https://www.pvv.ntnu.no/~stm/dring/driftflux-well.pdf}{preprint}). No Saurel-type seven-equation CO₂ decompression
validation was located; conditional hyperbolicity of single-pressure
six-equation models and the Baer--Nunziato remedy are standard knowledge
\textbf{{[}unverified primary{]}}.

\begin{longtable}[]{@{}
  >{\raggedright\arraybackslash}p{(\columnwidth - 10\tabcolsep) * \real{0.1667}}
  >{\raggedright\arraybackslash}p{(\columnwidth - 10\tabcolsep) * \real{0.1667}}
  >{\raggedright\arraybackslash}p{(\columnwidth - 10\tabcolsep) * \real{0.1667}}
  >{\raggedright\arraybackslash}p{(\columnwidth - 10\tabcolsep) * \real{0.1667}}
  >{\raggedright\arraybackslash}p{(\columnwidth - 10\tabcolsep) * \real{0.1667}}
  >{\raggedright\arraybackslash}p{(\columnwidth - 10\tabcolsep) * \real{0.1667}}@{}}
\toprule\noalign{}
\begin{minipage}[b]{\linewidth}\raggedright
Criterion
\end{minipage} & \begin{minipage}[b]{\linewidth}\raggedright
HEM (3 eq.)
\end{minipage} & \begin{minipage}[b]{\linewidth}\raggedright
HRM/HRM* (4 eq.)
\end{minipage} & \begin{minipage}[b]{\linewidth}\raggedright
HFM (4 eq. + nucleation)
\end{minipage} & \begin{minipage}[b]{\linewidth}\raggedright
Drift-flux (3--4 eq. + slip law)
\end{minipage} & \begin{minipage}[b]{\linewidth}\raggedright
Two-fluid (6--7 eq.)
\end{minipage} \\
\midrule\noalign{}
\endhead
\bottomrule\noalign{}
\endlastfoot
PDEs added to current solver & 0 & +1 (vapor mass) & +2--3 (bubble
number/area) & 0--1 (algebraic slip in fluxes) & +3--4 \\
Closures beyond EOS & none & relaxation time θ & nucleation rates,
growth law & C₀(α), v\_drift (\citealp{shi2005}/T2Well) & interfacial friction,
Δp regularization, per-phase wall laws \\
Closure uncertainty for CO₂ & none (by construction) & θ spans 0.04--2.8
ms, \textasciitilde60×; \textbf{condensation-side θ uncharacterized} &
heterogeneous rates ``not predictive'' & untested for CO₂ at low rates;
no-slip beat it at 6.5 MPa & ``would require local data that do not
exist'' (SINTEF) \\
Documented CO₂ fidelity & wave speeds good; plateau over-, discharge
under-predicted (0--26\% at chokes) & plateau corrected; MAPE 7.6\% &
MAPE 6.0\%, one parameter set & mispredicted FALCON holdup; needed for
fallback regime & did not clearly outperform HEM \\
Numerical robustness & sound-speed discontinuity at dome; otherwise
hyperbolic, conservative & smooth (frozen) acoustics --- \emph{better}
than HEM & stiff sources (implicit Euler) & preserves conservative
structure & conditional hyperbolicity (6-eq.); \textasciitilde2× state
size (7-eq.) \\
Relative cost & 1 (flash-dominated) & \textasciitilde1.1 &
\textasciitilde1.3--1.5 & \textasciitilde1.05 & ≥2 \\
Validation data available & SINTEF/Botros/COOLTRANS/Guo & same (plateau
level) & same & FALCON loop only (44 mm, 6.5 MPa) & none that
discriminates \\
\end{longtable}

\hypertarget{verdict-hem-is-defensible-as-the-starting-point-with-the-failure-modes-named-and-bounded}{%
\subsection{3. Verdict: HEM is defensible as the starting point --- with
the failure modes named and
bounded}\label{verdict-hem-is-defensible-as-the-starting-point-with-the-failure-modes-named-and-bounded}}

The working hypothesis survives critical testing on every quantitative
axis examined, provided its documented biases are treated as known
model-form error rather than discovered later.

\textbf{Time scales.} Fitted CO₂ flashing relaxation times are
0.04--2.82 ms (\citealp{log2024hrm}, \href{https://www.pvv.ntnu.no/~stm/dring/hfm_preprint.pdf}{HFM preprint}); the well drains through a 12.7 mm choke over minutes. Taking
θ ≤ 3 ms against a transport residence time of 10¹--10² s (mixture
velocities of order 1--10 m/s over 610 m), the equilibrium Damköhler
number is ≥10⁴--10⁵: local equilibrium is an excellent bulk
approximation everywhere except within the first metres behind a fast
rarefaction head and inside the choke itself (transit ≤
\textasciitilde0.25 ms at 50--100 m/s through a plate ≤ one diameter
thick --- firmly inside Fauske's \textasciitilde0.1 m non-equilibrium
regime, \href{https://public.lanl.gov/davidp/MSThesisHTML/2.html}{LANL
review}).

\textbf{Metastability.} The documented HEM plateau errors are for
dense/liquid inventories. For a gas-phase start, the relevant
non-equilibrium is delayed condensation --- observed once (droplets at
2.8 MPa vs higher equilibrium dew pressure, Test 3, \citealp{munkejord2020}, \href{https://www.pvv.ntnu.no/~stm/dring/depress-co2-exp.pdf}{preprint}) and bounded by CNT in \citet{hammer2025}. Near-critical dome entries show essentially zero
undershoot, and the well starts at reduced temperature ≈ 0.99--1.06, so
the warm early inventory crosses the dome nearly in equilibrium;
supercooling grows as the column cools. Crucially, \textbf{no measured
condensation-side relaxation time exists} --- the 0.04--2.8 ms values
are all flashing-side. This is simultaneously the weakest link in the
HEM justification and a publishable gap.

\textbf{Slip.} At 6.5 MPa in a 44 mm vertical pipe, no-slip beat every
drift-flux law tested
(\citealp{hammer2021}, \href{https://www.pvv.ntnu.no/~stm/dring/driftflux-well.pdf}{preprint}) --- but their own caveat, ``except at low inlet liquid
holdup,'' flags the regime where the blowdown problem lives. At 1--4 MPa
the liquid/vapor density ratio reaches order 10--50; the choked mass
flux through 12.7 mm, spread over a bore with \textasciitilde200× the
area, gives in-bore superficial velocities that fall to and below the
O(0.3--0.7 m/s) Kutateladze drift velocity. Below the Turner-type
critical gas velocity, condensate cannot be lifted: fallback, a liquid
sump on the closed bottom, and countercurrent segregation follow --- and
HEM then mispredicts inventory distribution and the DTS-visible
phase-front locations. This is an inference from cited density/geometry
facts plus standard flooding physics; no CO₂-well demonstration exists.

\textbf{Known mispredictions to expect.} Wave propagation: equilibrium
kink too sharp; measured transition smooth, between equilibrium and
frozen curves. Quality/dry-out: \textasciitilde0.5--2 s early in 61.7 m
pipes. Temperature: tracks T\_sat(p) correctly while two-phase (the
saturation clamp is HEM's most reliable output); the large temperature
errors observed are heat-transfer-closure failures, not equilibrium
failures. Discharge: 0--26\% underprediction at short restrictions,
which couples back into the pressure history (HEM at the choke
overpredicted the pipe plateau in \citealp{hammer2022}).

\textbf{Transitions and limits.} The single↔two-phase transition is
thermodynamically trivial in HEM (the flash decides) but numerically
delicate (Section 9). Near-critical: the lever rule degenerates and the
flash ill-conditions as ρ′−ρ″→0; published codes handle it by dense
tabulation, nothing smarter was found. Triple point (0.518 MPa, 216.59
K): reachable only late and essentially at/downstream of the choke; the
pure-CO₂ equilibrium sound speed is \textbf{zero at the triple point}
(\citealp{munkejord2021}, \href{https://www.pvv.ntnu.no/~stm/dring/depress-mix-exp.pdf}{preprint}), so the wellbore model should not carry a solid phase
--- a three-phase flash belongs in the choke submodel if at all (\citealp{hammer2013} provides the thermodynamic machinery).

\textbf{Global vs local reliability --- the honest summary.} HEM is
trustworthy for globally conserved and saturation-clamped quantities:
pressure histories, inventory, two-phase column temperature, wave
arrival to within the smoothing noted. It is untrustworthy for local,
closure-sensitive quantities: discharge rate through short restrictions,
local quality/void during fast transients, dry-out timing, phase
distribution at low rates, and anything downstream of the triple point.
The staged plan should be read as: Stage 1 earns the global quantities;
the later stages chase the local ones.

\hypertarget{hem-governing-equations-the-solver-changes-less-than-the-closures-do}{%
\subsection{4. HEM governing equations: the solver changes less than the
closures
do}\label{hem-governing-equations-the-solver-changes-less-than-the-closures-do}}

In conservative mixture form, with mixture density ρ, velocity u, total
energy E = ρ(e + u²/2), well inclination giving gravity component g\_z,
area A constant:

\[\frac{\partial \rho}{\partial t} + \frac{\partial (\rho u)}{\partial x} = 0\]

\[\frac{\partial (\rho u)}{\partial t} + \frac{\partial (\rho u^2 + p)}{\partial x} = -\rho g_z - \frac{f_D}{2D}\rho u |u|\]

\[\frac{\partial E}{\partial t} + \frac{\partial\big[(E + p)u\big]}{\partial x} = -\rho g_z u + \frac{4 q''_w}{D}\]

These are \emph{identical in form} to the existing single-phase system.
Under full equilibrium, ρ and e are mixture quantities built from
saturated phase states at a common (p, T\_sat(p)):

\[\frac{1}{\rho} = v = (1-x)\,v'(T) + x\,v''(T), \qquad e = (1-x)\,e'(T) + x\,e''(T),\]

with quality x the lever-rule coordinate. Unchanged: the finite-volume
update, MUSCL-Hancock reconstruction (reconstruct in primitive variables
(ρ, u, p) as now), the conjugate completion/formation radial network and
its coupling through q″\_w, the gravity and geometry treatment, and the
fiber-optic output machinery. Changed: (i) the thermodynamic closure
(ρ,e)→(p,T,a) must return two-phase states, quality, and the
\textbf{equilibrium} sound speed a\_eq inside the dome (Section 5); (ii)
the friction closure needs a two-phase evaluation (Section 7); (iii) the
wall heat-transfer coefficient needs regime logic (Section 8); (iv) the
choke boundary needs a two-phase critical-flux routine and, per the
evidence, a non-equilibrium correction (Section 10). No new PDE, no new
wave family: eigenvalues remain u, u ± a with a = a\_eq(ρ,e) in the
dome. The full-equilibrium sound speed is (Flåtten \& Lund,
\href{https://www.sintef.no/globalassets/project/co2-dynamics/publications/subchar.pdf}{subchar.pdf}):

\[\frac{1}{a_{\mathrm{HEM}}^2} = \frac{1}{a_p^2} + \rho T\!\left[\frac{\rho_g \alpha_g}{c_{p,g}}\left(\frac{ds_g}{dp}\right)_{\!\mathrm{sat}}^2 + \frac{\rho_l \alpha_l}{c_{p,l}}\left(\frac{ds_l}{dp}\right)_{\!\mathrm{sat}}^2\right], \qquad \frac{1}{a_p^2} = \rho\left(\frac{\alpha_g}{\rho_g c_g^2} + \frac{\alpha_l}{\rho_l c_l^2}\right),\]

where a\_p is the Wood (mechanical-equilibrium) speed and the
saturation-line entropy derivatives supply the extra compliance from
phase change. Note a\_HEM ≤ a\_Wood ≤ a\_frozen --- the practical
consequence is that the explicit CFL constraint inside the dome is
\emph{looser} than in single phase; the time step stays governed by the
single-phase region.

\hypertarget{thermodynamic-closure-the-ux3c1e-flash-reduces-to-one-scalar-equation-in-t}{%
\subsection{5. Thermodynamic closure: the (ρ,e) flash reduces to one
scalar equation in
T}\label{thermodynamic-closure-the-ux3c1e-flash-reduces-to-one-scalar-equation-in-t}}

The published pattern (\citealp{giljarhus2012}, the foundational density--energy Span--Wagner method for CO₂ CFD; extended to vapor--liquid--solid in \citealp{hammer2013} ---
\href{https://pvv.ntnu.no/~stm/dring/publ.html}{Munkejord publ. list})
maps directly onto the existing code:

\textbf{Two-phase detection.} Given a cell's (ρ, e): compute the
saturated bracket energies e′\_sat(ρ), e″\_sat(ρ) from a precomputed 1-D
saturation-line table (T\_sat, p\_sat, ρ′, ρ″, e′, e″, s′, s″
parameterized by T). If e lies between the brackets (and ρ between
ρ″(T\_tr) and ρ′ limits), the state is in the dome; otherwise the
existing single-phase (ρ,e)→T Newton applies unchanged.

\textbf{The reduced flash.} For a pure substance, the general UV-flash
is 4 nonlinear equations, but saturation relations collapse it to
\textbf{one scalar equation in T}
(\citealp{kumar2024}, \href{https://arxiv.org/html/2408.09164v2}{arXiv:2408.09164}): with v = 1/ρ fixed,

\[F(T) = \big(1 - x(T)\big)\,e'(T) + x(T)\,e''(T) - e = 0, \qquad x(T) = \frac{v - v'(T)}{v''(T) - v'(T)},\]

solved by Newton with bisection safeguard inside the bracket; then p =
p\_sat(T), and a = a\_eq from the formula above (or the equivalent
(∂p/∂ρ)\_s with equilibrium-constrained derivatives, as SINTEF does).
\citet{kumar2024} report an order-of-magnitude CPU reduction on a tank
blowdown and 3--4× on an HLLC pipe case from this reduction alone. The
lever rule applies on v and e simultaneously by construction;
\textbf{quality x (mass fraction) and void fraction α = x v″/v must
never be conflated} --- α is what gravity head and DTS-visible density
respond to, x is what the energy balance transports; at 2 MPa, x = 0.1
corresponds to α ≈ 0.7--0.8.

\textbf{Acceleration.} Direct Span--Wagner evaluation can consume
\textasciitilde90\% of runtime; SINTEF's (T,p)-indexed look-up tables
achieved \textbf{100--300× speedups at \textasciitilde10⁻³ \% deviation}
(\href{https://www.sintef.no/en/publications/publication/1274087}{SINTEF
2007 record}). For the dome, the existing 512² (ln ρ, T) tables are
\emph{not} usable, since (ρ,T) is not a complete state pair under
saturation. Two published options: (a) \textbf{separate 1-D
saturation-line splines + reduced flash} --- the minimal change,
recommended; (b) a full (ρ,e)-space 2-D table spanning the dome ---
Brown et al.~(UCL) built adaptive inverse-interpolation tables on a
(ρ,U) grid with \textbf{70\% of nodes inside the envelope, 15\% each
hugging the dew and bubble lines}, Akima splines, mean deviations
0.005\% (T) and 0.44\% (p); a uniform (p,T) grid ``failed outright''
(\href{https://discovery.ucl.ac.uk/1522588/1/Brown\%20et\%20al\%20Thermodynamic\%20interpolation\%20for\%20the\%20simulation\%20of\%20two-phase\%20flow\%20of\%20non-ideal\%20mixtures\%20AAM.pdf}{UCL
Discovery PDF}). A rumored standalone ``SINTEF SW-spline'' journal paper
\textbf{does not verifiably exist} --- do not cite one.

\textbf{Near-critical behavior.} The lever rule degenerates as ρ′−ρ″ → 0
and the saturation-line Newton ill-conditions; no published special fix
was found beyond dense tabulation near T\_c --- flag this as an
implementation risk to be handled by local grid refinement of the
saturation table and a fallback to the single-phase branch within a
small ε of T\_c.

\textbf{CoolProp organization (background orientation, not from the
research notes --- standard knowledge, verify against CoolProp
documentation before citing):} CoolProp implements the Span--Wagner
Helmholtz formulation with (ρ,T) as native variables, uses ancillary
saturation equations as initial guesses for a Newton saturation solve,
routes general two-property flashes (including (ρ,e)\ldots{} via
\texttt{DmassUmass}) through 1-D/2-D Newton iterations with phase
detection against the saturation envelope, and offers tabular backends
(TTSE and bicubic) for order-of-magnitude speedups --- i.e., the same
architecture as options (a)/(b) above, and a useful reference
implementation and cross-check for the in-house flash.

\hypertarget{mixture-transport-properties-effective-properties-are-regressions-not-physics}{%
\subsection{6. Mixture transport properties: effective properties are
regressions, not
physics}\label{mixture-transport-properties-effective-properties-are-regressions-not-physics}}

The critical conceptual distinction, which the report's downstream
choices all rest on: a \textbf{constitutive property} (μ, k of a single
phase) relates local flux to local gradient and is measurable
independent of geometry; an \textbf{effective mixture property} (μ\_m,
k\_m) is a homogenization artifact defined only relative to the model
that consumes it --- ``whatever makes f(Re\_m) reproduce Δp''; a
\textbf{two-phase multiplier} (φ²\_l: Lockhart--Martinelli, Friedel,
Müller-Steinhagen--Heck) abandons the effective-fluid fiction and
correlates the pressure-gradient ratio directly --- empirically superior
precisely because it does not pretend a continuum property exists; a
\textbf{wall-HTC correlation} is a boundary-condition closure, not a
property at all. Error chains must not cross categories: a μ\_m tuned
for friction fixes nothing about Pr\_m for heat transfer.

\textbf{Viscosity.} Verified forms and years
(\href{https://aghajar.okstate.edu/sites/default/files/pub-pdfs/Ghajar\%20\%26\%20Bhagwat_HTE\%2034(13)_pp\%201044-1059\%20(2013).pdf}{Ghajar \& Bhagwat 2013 PDF}; \citealp{ghajar2013}): \citet{mcadams1942} 1/μ\_m = x/μ\_g + (1−x)/μ\_l;
\citet{cicchitti1960} μ\_m = xμ\_g + (1−x)μ\_l; \citet{dukler1964} volume-weighted; \citet{beattie1982} Einstein-suspension hybrid on
homogeneous void; \citet{awad2008} Maxwell-bound
effective-property forms. Their quantitative verdict on 616 annular
points across 9.5--50 mm pipes: the best model (Awad--Muzychka 1)
captured only \textbf{60.4\% of data within ±30\%}, Cicchitti 57.6\% ---
and ``none \ldots{} predict accurately for all pipe sizes.'' The
Thermopedia assessment is blunter: the homogeneous model works only at
high velocities with strongly interdispersed phases, and ``adjusting the
viscosity definition alone does not fix this''
(\href{https://www.thermopedia.com/de/content/37/}{Thermopedia}). Two
mitigating facts for this application: μ\_m enters turbulent friction
only as μ\_m\^{}0.25, so a factor-2 disagreement moves Δp\_f by
\textasciitilde19\%; and as p → p\_c all definitions converge (μ\_l/μ\_g
\textasciitilde{} 3--10 for near-critical CO₂). Recommendation:
\textbf{McAdams (or Cicchitti) baseline, Beattie--Whalley as the
sensitivity alternative, and an explicit statement that at G = 10--100
kg/m²s in a 170 mm bore the gravity head --- i.e., the void fraction ---
dominates the pressure gradient, not friction.} Convenience averages for
μ\_m are tolerable only because the quantity barely matters; they must
not be promoted into k\_m or c\_p,m by analogy.

\textbf{Conductivity.} No verified source endorses an effective
two-phase k\_m as a wall-transfer property, and verified practice is
uniformly against it: the UCL CO₂ blowdown code (Dittus--Boelter +
Rohsenow ---
\href{https://discovery-pp.ucl.ac.uk/id/eprint/10064608/1/matecconf_icchmt2018_01026.pdf}{UCL
MATEC 2018}), the Cheng--Ribatski--Thome CO₂ model
(\href{https://thermopedia.com/content/10372/}{Thermopedia}), and \citet{shah2017} all build h on single-phase properties plus two-phase
enhancement/suppression. The physical reason: wall transfer is
controlled by near-wall structure (liquid film, nucleation, dry
patches); the wall sees k\_l or k\_g, never their average. k\_m is
legitimate only as a stagnant-dispersion conduction property
(Wiener/Maxwell bounds) --- relevant solely in the late-blowdown
near-stagnant limit --- or inside a deliberately crude ``effective
Dittus--Boelter'' sensitivity bracket, whose error no published study
has even quantified (a stated negative result).

\hypertarget{wall-friction-homogeneous-darcyweisbach-baseline-muxfcller-steinhagenheck-as-the-sensitivity-upgrade}{%
\subsection{7. Wall friction: homogeneous Darcy--Weisbach baseline,
Müller-Steinhagen--Heck as the sensitivity
upgrade}\label{wall-friction-homogeneous-darcyweisbach-baseline-muxfcller-steinhagenheck-as-the-sensitivity-upgrade}}

The HEM momentum equation's self-consistent closure is homogeneous
Darcy--Weisbach with f\_D from the existing Churchill correlation
evaluated at Re\_m = GD/μ\_m --- retaining the current friction
architecture verbatim. The best-supported simple alternative is
\textbf{\citet{muller1986}}: ranked first (with Grönnerud)
on EPFL refrigerant-evaporation data, with Friedel third at
\textasciitilde30\% standard deviation single-component
(\href{https://iifiir.org/en/fridoc/prediction-of-two-phase-pressure-gradients-of-refrigerants-in-119317}{Ould Didi, Kattan \& Thome 2002 record}; \citealp{oulddidi2002};
\href{https://www.thermopedia.com/de/content/37/}{Thermopedia}). Its
structure --- dp/dx\textbar\_f = {[}A + 2(B−A)x{]}(1−x)\^{}\{1/3\} + Bx³ with A, B the all-liquid and all-gas single-phase gradients
\textbf{{[}formula from training, 1986 paper not fetched{]}} --- is
C⁰-continuous in x with exact single-phase limits, which matters for
transient stability. The famous databank confirmations (Tribbe \&
Müller-Steinhagen 2000; Xu \& Fang 2012, MAD ≈ 28.8\% often quoted) were
\textbf{not verified this session --- do not quote their statistics}.

Honest applicability statement: every databank found is D ≲ 50 mm and G
≳ 100 kg/m²s; \textbf{no large-databank friction study at D
\textgreater{} 100 mm exists at all} --- a genuine literature gap. All
steady correlations are applied quasi-steadily per cell per step; the
defense is time-scale separation (friction response \textasciitilde D/u*
≪ blowdown time), which is an assumption, not a result --- no source
quantifies the quasi-steady error. Recommended posture: homogeneous
f(Re\_m(μ\_McAdams)) baseline; MSH as the sensitivity alternative;
report friction-closure uncertainty as subordinate to
void-fraction/gravity-head uncertainty, which is structural to HEM.

\hypertarget{wall-heat-transfer-regime-logic-under-conjugate-coupling-with-every-correlation-an-extrapolation-in-diameter}{%
\subsection{8. Wall heat transfer: regime logic under conjugate
coupling, with every correlation an extrapolation in
diameter}\label{wall-heat-transfer-regime-logic-under-conjugate-coupling-with-every-correlation-an-extrapolation-in-diameter}}

The conceptual anchor: in blowdown, the latent-heat sink is the
\textbf{bulk} --- depressurization flashes/condenses the mixture and the
fluid temperature tracks T\_sat(p) regardless of the wall --- while the
wall HTC sets only the conjugate heat leak from completion and formation
into the fluid. Wall \textbf{nucleate boiling} is a separate phenomenon
that switches on only when T\_w − T\_sat exceeds the onset superheat.
The direct precedent (Teng/UCL, HEM blowdown of a 50 mm steel pipe with
lumped conjugate wall) used exactly this logic --- Dittus--Boelter
single-phase, Rohsenow nucleate boiling two-phase --- and produced the
application-critical verified finding: \textbf{``the rate of CO₂ solid
formation in the pipe is limited by heat transfer at the pipe wall,''}
with wall heating able to evaporate liquid before the triple point is
crossed, and adiabatic treatment overestimating solid mass by 15--50\%
(\href{https://discovery-pp.ucl.ac.uk/id/eprint/10064608/1/matecconf_icchmt2018_01026.pdf}{UCL
MATEC 2018}). For a finite-inventory well with a conjugate network
already built, wall heat input sets the duration of the
saturation-clamped period --- Lindeberg's adiabatic conclusion does not
apply.

Correlation inventory vs the application (170 mm, G \textasciitilde{}
10--100 kg/m²s, 220--320 K): the CO₂-specific
Cheng--Ribatski--(Wojtan)--Thome model is validated for \textbf{D =
0.529--9.52 mm and G ≥ 100 kg/m²s}
(\href{https://thermopedia.com/content/10372/}{Thermopedia}) --- 18×
under the bore diameter with a mass-flux floor at the application's
ceiling: formally inapplicable; its portable piece is the
reduced-pressure nucleate-boiling scaling. \textbf{\citet{shah2017}} is the least-bad extrapolation: 4852
points, 30 fluids including CO₂, D = 0.38--27.1 mm, G = 15--2437 kg/m²s,
reduced pressure to 0.787, \textbf{MAD 18.6\%}
(\href{https://iifiir.org/en/fridoc/unified-correlation-for-heat-transfer-during-boiling-in-plain-140363}{IIR
record}) --- still a 6× diameter extrapolation. \textbf{No CO₂
flow-boiling data above \textasciitilde10 mm diameter were located;
every two-phase wall-HTC evaluation in this well is an extrapolation in
D and mostly in G} --- frame the HTC as a sensitivity study (±50\%
blanket), not a validated closure.

Recommended initial implementation, mirroring verified practice and
LOCA-code structure (RELAP5 mode-ladder specifics are \textbf{{[}T ---
manuals located at NRC ADAMS, not fetched{]}}): (i) single-phase ---
Dittus--Boelter/Gnielinski with a \textbf{natural-convection floor}
(Churchill--Chu-type), because h ∝ G\^{}0.8 → 0 unphysically as flow
dies while buoyant convection persists in a 170 mm vertical bore ---
this is precisely the failure the SINTEF N₂ closed-end test exposed
(−150 °C predicted vs −78 °C measured); (ii) two-phase with T\_w
\textgreater{} T\_sat --- Chen-type superposition or \citet{shah2017}, plus a
pool-boiling floor (Rohsenow per the UCL precedent, or Cooper
reduced-pressure form \textbf{{[}Cooper 1984 unverified{]}}); prefer
smooth superposition/asymptotic blends (Chen S-factor; Liu--Winterton h²
form \textbf{{[}unverified{]}}) over hard if/else switches, which inject
HTC discontinuities into the conjugate wall ODEs; (iii) two-phase with
T\_w \textless{} T\_sat (condensation on cold walls, the relevant mode
when warm wellbore gas contacts deep cold completions or during late
repressurization) --- Shah condensation \textbf{{[}1979/2009 years
unverified{]}}; (iv) h = max(forced, natural) within branches for regime
continuity. Stated negative result: flashing-specific wall-HTC
literature at near-zero wall superheat is essentially absent; the
low-velocity bottomhole regime (near-stagnant two-phase natural
convection in a large vertical enclosure) has no published closure at
all --- the weakest closure in the model, by construction.

\hypertarget{numerical-implications-keep-hllc-respect-the-dome-boundary-kink}{%
\subsection{9. Numerical implications: keep HLLC, respect the
dome-boundary
kink}\label{numerical-implications-keep-hllc-respect-the-dome-boundary-kink}}

\textbf{Hyperbolicity is not the problem.} The sum-of-squares
construction guarantees a real a\_eq for any thermodynamically stable
EOS, so the HEM system is unconditionally hyperbolic
(\href{https://www.sintef.no/globalassets/project/co2-dynamics/publications/subchar.pdf}{Flåtten
\& Lund, subchar.pdf}). The pathology is loss of \emph{smoothness}:
p(ρ,e) has a kink at the dome boundary and the sound speed jumps (CO₂ at
5 MPa, 287.43 K: c\_g = 201.5 m/s, c\_l = 398.9 m/s single-phase, versus
a two-phase equilibrium value of order tens of m/s --- the decompression
plateau at 6 MPa computes to \textasciitilde40 m/s; the water analog
drops 1543 → \textasciitilde1 m/s on the liquid side). Only full p--T--μ
equilibrium produces the discontinuity --- HRM's frozen acoustics are
smooth, which is why the HEM→HRM stage \emph{improves numerics as well as physics} (\citealp{lund2012}, \href{https://www.sintef.no/globalassets/project/co2-dynamics/publications/lund_two-phase_relaxation_hierarchy.pdf}{hierarchy PDF}).

\textbf{Expected wave anatomy.} The kink makes p(ρ,e) behave like a
Menikoff--Plohr piecewise-smooth EOS: rarefactions crossing the dome
\textbf{split into composite double fans}, directly observed in CO₂ HEM Test 4 of \citet{log2021}, and captured by their HLLC-type solvers as a kinked fan
without special treatment, at locally first-order accuracy
(\href{https://www.pvv.ntnu.no/~stm/dring/hllcs-hem.pdf}{preprint}).
\citet{kumar2024} attribute their degraded convergence orders to the
discontinuous flash derivatives, not instability
(\href{https://arxiv.org/html/2408.09164v2}{arXiv:2408.09164}).

\textbf{Published HLLC experience for HEM CO₂.} No examined source
reports HLLC failing \emph{at the dome boundary} in a straight
conservative duct; the documented failures (negative internal energy,
spikes) in \citet{log2021} are driven by the non-conservative area-change
term and resonance --- absent from this constant-area well. \citet{kumar2024} run plain HLLC + forward Euler on CO₂ HEM with no reported
oscillations. SINTEF's production choice of the more dissipative FORCE
flux at CFL 0.85 for its validation code reads as robustness-first
pragmatism, and signals that extra dissipation at mixed faces is an
accepted trade. Brown et al.~(UCL) ran AUSM+ at CFL 0.2 stably on
two-phase shock tubes.

\textbf{Recommended strategy} (inference from adjacent evidence; no
systematic HLLC-vs-fallback failure statistics exist --- stated gap):
(1) retain HLLC with Davis/Einfeldt wave-speed estimates, taking at
faces that straddle the dome the \textbf{maximum of the left/right sound
speeds} --- conservative by the subcharacteristic ordering, so no
equilibrium wave can outrun the estimate; (2) guard positivity by
checking intermediate-state density and internal energy and locally
falling back to HLL/Rusanov on violation (\citealp{log2021}: pressure-based intermediate variables help; negative internal energy was the observed
failure mode); (3) flash every cell after each stage with phase
detection against the saturation brackets, flooring x ∈ {[}0,1{]} and
clipping tiny negative vapor masses; (4) accept locally first-order
resolution of the dome-crossing kink; refine the grid rather than
engineer the flux; (5) note the CFL constraint inside the dome is
\emph{looser} (a\_HEM \textless{} a\_single-phase) --- the time step
remains governed by single-phase regions; the late-blowdown hazards are
flash non-convergence at low pressure and strong-rarefaction positivity,
handled in published codes by robust flash fallbacks and dissipative
fluxes rather than formal positivity-preserving limiters (none of the
examined CO₂ papers uses one --- stated negative result).

\hypertarget{the-choke-equilibrium-wellbore-plus-non-equilibrium-restriction-is-the-published-physically-justified-hybrid}{%
\subsection{10. The choke: equilibrium wellbore plus non-equilibrium
restriction is the published, physically justified
hybrid}\label{the-choke-equilibrium-wellbore-plus-non-equilibrium-restriction-is-the-published-physically-justified-hybrid}}

The residence-time split is decisive: the 610 m wellbore is
\textasciitilde6,000 Fauske relaxation lengths (local equilibrium
overwhelmingly justified); the ≤12.7 mm-thick plate is a small fraction
of one (transit ≤ \textasciitilde0.25 ms, far below the \textasciitilde1
ms-scale flashing relaxation) --- firmly non-equilibrium. \textbf{\citet{hammer2022} did exactly this}: an equilibrium HEM transient pipe
solver with Henry--Fauske and delayed-HEM restriction models as the
outflow boundary, finding the hybrid matches CO₂ blowdown data better
than HEM-everywhere, with HF best for the pipe cases
(\href{https://www.pvv.ntnu.no/~stm/dring/outflow_preprint.pdf}{preprint}).
The coupling pattern is also the current code's: evaluate the
quasi-steady restriction model from the boundary-cell state treated as
stagnation, area-scale the returned mass flux, impose it at the
boundary; declare choking when velocity reaches the local sound speed
along the expansion path, else set the flux by back-pressure. Two
cautions from the same source: the boundary-cell-as-stagnation
approximation neglects the pipe-side velocity head (second-order except
late in blowdown), and the HEM G-maximization has a kink at the
saturation pressure --- \textbf{bracket the saturation point explicitly
and search the single-phase and two-phase branches separately}, or
Newton searches oscillate; HEM ``often'' chokes exactly at the
saturation crossing.

\textbf{The C\_d tension must be resolved, not papered over.} The
physically grounded contraction coefficient for a sharp-edged orifice
passing flashing CO₂ is \textbf{C\_c ≈ 0.74--0.75} (\citealp{hammer2022}, Bragg-type force balance from the incompressible π/(π+2) ≈ 0.611 limit),
while the current C\_d = 0.84 sits in the compressible-gas-orifice
range. The two errors currently compensate: HEM underpredicts G by
\textasciitilde0--26\%, and the high C\_d pulls the product back toward
the data. This is fortuitous, state-dependent, and must be made
explicit: \textbf{switching to HF or HNE-DS while keeping C\_d = 0.84
would overpredict.} Recommended: for the gas-dominated early blowdown,
HEM flux with C\_d = 0.84 is defensible; once the choke-inlet state is
two-phase or liquid-like, move to (HF or D-HEM) paired with C\_c ≈ 0.75,
and calibrate the (model, C\_d) pair jointly against the facility's own
Coriolis/rate data. If the device is a choke bean (L/D \textasciitilde{}
2--6) rather than a thin plate, partial reattachment raises effective
C\_d toward 0.8--0.9 --- geometry must be pinned down before
calibration. Interface consistency: the imposed boundary flux must carry
mass and total energy (including kinetic and gravitational terms)
consistently with the interior conservative update; the hybrid's
thermodynamic inconsistency (the choke generates entropy/metastability
the wellbore model forbids) is confined to the boundary cell and is the same approximation \citet{hammer2022} accepted --- but beware
double-counting: if the wellbore HEM delivers a boundary state whose
quality is an equilibrium artifact, HF's non-equilibrium correction
partially re-corrects vapor that was never really there; this matters
most when the flash front sits near the wellhead. Choked/unchoked
switching must be continuous: as upstream pressure decays, G\_choked
must meet the subsonic orifice flux at the matching back-pressure; the
HEM sound-speed discontinuity can make this switch non-smooth and needs
regularization in code. No formal error analysis of the
equilibrium-pipe/non-equilibrium-choke split exists (practiced, not
proven --- stated gap).

\hypertarget{the-staged-hierarchy-is-correctly-ordered-with-one-re-prioritization}{%
\subsection{11. The staged hierarchy is correctly ordered --- with one
re-prioritization}\label{the-staged-hierarchy-is-correctly-ordered-with-one-re-prioritization}}

The proposed ladder HEM → improved transport closures → HRM →
drift-flux/two-fluid is broadly right; the critical assessment changes
one thing: \textbf{the non-equilibrium choke correction (part of the HEM
stage's boundary condition) and the HRM stage outrank transport-closure
refinement}, because the documented first-order errors (discharge rate,
plateau) are phase-kinetic, while friction/HTC refinements move
second-order terms in this G--D regime. The gravity-head/void error, by
contrast, is structural and is what ultimately forces the drift-flux
stage.

\textbf{Level 0 --- HEM (implement now).} Added equations: none.
Closures: equilibrium flash, McAdams μ\_m + Churchill, regime-logic HTC,
HEM choke flux with explicit saturation-bracketing. Gains: full
two-phase capability, saturation-clamped temperatures, correct inventory
bookkeeping. Cost: flash-dominated; mitigate with saturation-line
splines. Numerical risk: dome-boundary kink (Section 9). Validation: SINTEF Test 3 + dense tests, Botros curves \citep{munkejord2020,botros2016}.

\textbf{Level 1 --- transport-closure upgrades + non-equilibrium choke
(sensitivity).} Added equations: none. Closures: MSH friction, \citet{shah2017}/Chen HTC with floors, HF or D-HEM at the choke with C\_c ≈ 0.75.
Gains: bounds the closure sensitivity; fixes the one documented
first-order HEM bias (discharge). Cost: trivial. Validation: \citet{hammer2022} orifice data (exact 12.7 mm size), COOLTRANS vents.

\textbf{Level 2 --- HRM* (implement when evidence demands).} Added
equations: \textbf{one} --- vapor-mass with Γ = (ρy\_eq − ρy)/θ,
integrated implicitly (Godunov splitting per \citealp{log2024hfm}). Closures: θ,
treated as an uncertainty parameter spanning 0.05--3 ms (a single
constant does not transfer across initial states --- 60× spread; a
Downar-Zapolski-form state dependence refit to CO₂ follows the
Angielczyk/Haida line, originals \textbf{{[}unverified{]}}). Properties:
metastable single-phase extensions of Span--Wagner. Gains: smooth
(frozen) acoustics at the dome --- a numerical improvement --- plus
corrected plateau/undershoot; the honest caveat is that \textbf{the
condensation-side θ for the gas-phase well is uncharacterized} --- the one dataset is \citet{hammer2025}, which should be
obtained in full. Cost: \textasciitilde10\%. Validation: \citet{log2024hrm} initial-state study.

\textbf{Level 3 --- drift-flux slip (defer until evidence).} Added
equations: none to the hyperbolic core if added algebraically in the
fluxes (one extra PDE if phasic mass is tracked). Closures: \citet{shi2005} C₀ and drift velocity as modified by \citet{pan2011} (Ku = 142, smoothed regime interpolation) --- take forms from the SPE
original. Gains: condensate fallback, sump formation, countercurrent
segregation --- the late-blowdown physics HEM structurally cannot
represent, and the regime where DTS phase-front locations go wrong.
Trigger quantitatively defined in Section 14. Validation: only the
FALCON loop (44 mm, 6.5 MPa, wrong regime) exists --- this level will be
validated by the PERTT experiment itself or not at all.

\textbf{Level 4 --- two-fluid (not recommended).} +3--4 PDEs,
interfacial closures that ``would require detailed local experimental
data that do not exist'' (SINTEF's own words), conditional hyperbolicity
management, \textasciitilde2× cost --- against a record in which the
six-equation model failed to beat HEM. Defensible only if the experiment
produces local void/slip data that drift-flux demonstrably cannot fit.

\hypertarget{verification-validation-and-uncertainty-a-composite-ladder-by-necessity}{%
\subsection{12. Verification, validation, and uncertainty: a composite
ladder by
necessity}\label{verification-validation-and-uncertainty-a-composite-ladder-by-necessity}}

\textbf{Numerical verification.} (1) Exact conservation of mass and
total energy to round-off over long closed-boundary runs, including
dome-crossing cells. (2) Stationarity: a hydrostatic, geothermal
two-phase column (lower liquid, upper vapor, interface at T\_sat(p))
must remain at rest --- the two-phase analog of the current
well-balancing test. (3) Single-phase limit: results identical to the
existing solver for trajectories that never enter the dome, and
continuous as a trajectory grazes the dew line. (4) Saturation-crossing
1-D Riemann problems: the 200 m tube of \citet{kumar2024} (100 bar/300 K ‖ 30
bar/300 K, t = 0.2 s) with published reference profiles, and \citet{log2021} Test 4 exhibiting the diagnostic \textbf{split rarefaction}
(\href{https://arxiv.org/html/2408.09164v2}{arXiv:2408.09164};
\href{https://www.pvv.ntnu.no/~stm/dring/hllcs-hem.pdf}{preprint}). (5) Two-phase acoustics: small-amplitude waves in a
uniform two-phase state propagating at the computed a\_HEM. (6)
Grid/time convergence: expect second order in smooth regions degrading
to first at the dome kink --- document, don't fight. (7) Positivity
stress tests: strong double-rarefactions into near-vacuum. (8)
\citet{edwards1970} horizontal pipe blowdown \textbf{{[}page range via secondary sources{]}} for
cross-community credibility --- note it is water, so it exercises the
flash machinery only if an IAPWS backend is added; optional.

\textbf{Validation datasets and what each validates.} \emph{SINTEF ECCSEL} \citep{munkejord2020,sintef2020data}:
Tier-1, open raw data, complete geometry/sensor/uncertainty
documentation. \textbf{Gas-phase Test 3 (4.04 MPa, 10.2 °C) is the
single most important case} --- nearly the well's state, exhibiting
decompression-induced condensation; validates wave speeds, condensation
onset, pressure histories. Dense Tests 4/6/8 validate the two-phase
closure and (with the N₂ reference) expose heat-transfer-closure failure
modes. \emph{CO₂--N₂/He companion} \citep{munkejord2021,sintef2021data} adds mixture robustness if ever needed.
\emph{\citet{botros2016} pure-CO₂ shock tube}:
decompression-wave speed vs pressure --- the cleanest
EOS/characteristics test. \emph{COOLTRANS vertical vents} (12/25/50 mm
at 150 barg with Coriolis mass flow): the surface-choke
boundary-condition benchmark, including the 50 mm delayed-nucleation
excess. \emph{Guo et al.~Test 1} (gas, 3.7 MPa, 310.6 K, vertical 15 mm
orifice, 258 m): large-scale warm gas-phase cross-check closest to the
well. \emph{\citet{hammer2022} orifice blowdowns}: the choke model at
the exact 12.7 mm size. \emph{Vessel blowdowns} (INERIS sphere once retrieved; \citealp{gebbeken1996} once verified): long-time energy balance,
flashing onset, stratification --- the analog of the well's slowly
draining lower inventory. \textbf{No dataset satisfies all of (vertical,
\textgreater100 m, gas-phase \textasciitilde5 MPa, choked release,
distributed axial P/T): validation is composite by construction, and the
PERTT blowdown would be the first dataset to close that gap.}

\textbf{Dominant model-form uncertainties per predicted quantity.}
Discharge rate: choke non-equilibrium model + C\_d (±10--25\%, the
largest single lever). Pressure history: choke flux (which sets drainage
rate) then EOS (negligible with Span--Wagner). Two-phase column
temperature: essentially pinned to T\_sat(p) --- small model-form error;
single-phase-region temperature: wall-HTC closure (demonstrated −72
°C-class errors at low velocity in the N₂ test). Phase-front
locations/void profile: no-slip assumption (unbounded at low rates ---
structural). Late-time behavior near the triple point: wall heat input
(solid formation is wall-HT-limited per the UCL precedent) and the
uncharacterized condensation kinetics. Propagate the choke (C\_d,
model), θ (0.05--3 ms), HTC (±50\%), and μ\_m/friction alternates
through a sensitivity matrix rather than a single calibrated run.

\hypertarget{dts-interpretation-the-saturation-thermometer-is-sound-physics-and-an-unclaimed-observation}{%
\subsection{13. DTS interpretation: the saturation thermometer is sound
physics and an unclaimed
observation}\label{dts-interpretation-the-saturation-thermometer-is-sound-physics-and-an-unclaimed-observation}}

In an equilibrium two-phase column, T = T\_sat(p(z)): a DTS trace plus a
single pressure gauge pins the pressure profile and locates phase fronts
--- the ``saturation thermometer.'' The thermodynamic basis is solid and
implicitly recognized at Ketzin (whose operators flushed with N₂
precisely because condensation/evaporation rate asymmetry causes
pressure problems ---
\href{https://gfzpublic.gfz-potsdam.de/pubman/item/item_242409_1/component/file_242408/15975.pdf}{Würdemann
et al.~2010}), \textbf{but no published DTS trace of a two-phase CO₂
column in a well exists} --- Ketzin's record is single-phase by design,
and the geothermal flash-point analog, while standard practice, was not
verified to a citable source this session. This is the clearest open
observational niche the planned experiment can claim.

Expected signatures, in model-output terms: (i)
\textbf{flashing/condensation onset} --- the depth at which the local T
departs from the single-phase expansion-cooling trend and locks onto
T\_sat(p(z)), a kink that propagates as the dome boundary moves; (ii)
\textbf{plateaus} --- along a two-phase interval the DTS trace follows
the saturation curve's gentle hydrostatic slope dT\_sat/dz =
(dT/dp)\_sat · ρ\_m g, qualitatively distinct from both the geothermal
gradient and the single-phase transient; (iii) \textbf{slip/segregation
signatures} --- HEM predicts a single connected two-phase interval,
whereas fallback and sump formation produce a liquid-rich bottom (T at
T\_sat but with the mixture-density-driven pressure profile of a liquid
column) and possibly a dry middle section; a saturation plateau whose
implied hydrostatic gradient disagrees with the no-slip ρ\_m(z) is
direct slip evidence; (iv) \textbf{post-shut-in evaporative cooling} at
the receding liquid surface. HEM-vs-HRM distinguishability on DTS is
marginal (θ \textasciitilde{} ms against DTS integration times of
seconds) \emph{except} at condensation onset, where supercooled vapor
would show T temporarily \textbf{below} T\_sat(p) --- a few-K, transient
signature that is the one field-observable discriminator of condensation
non-equilibrium, and precisely the uncharacterized quantity (Section 3).
Additional model outputs to emit now, cheaply: T\_sat(p(z,t)) alongside
T(z,t); the superheat/supercool margin T − T\_sat; quality x(z,t) and
void α(z,t) separately; the phase-boundary depths vs time; and synthetic
DTS traces with the instrument's spatial/temporal averaging applied, so
model-data comparison is apples-to-apples.

\hypertarget{recommendations-ah-and-the-prioritized-roadmap}{%
\subsection{14. Recommendations A--H and the prioritized
roadmap}\label{recommendations-ah-and-the-prioritized-roadmap}}

\textbf{A. HEM as the start? Yes.} Grounds: Damköhler ≥10⁴ (θ =
0.04--2.8 ms vs minutes), near-zero undershoot for warm dome entries,
no-slip winning on the only vertical CO₂ data, two-fluid failing to beat
HEM, and architecture preservation. Conditions: non-equilibrium
correction at the choke from the outset; dome-boundary numerics per
Section 9; the known biases logged as expected model-form error.

\textbf{B. Thermodynamic closure.} Span--Wagner (not cubic --- the
PR-vs-OLGA temperature bias is documented) via the reduced UV-flash:
saturation-line 1-D splines + one-scalar Newton in T with lever-rule
quality on v and e; keep the existing single-phase tables and Newton
untouched outside the dome; consider a Brown-style (ρ,e) dome table
(dense dew/bubble-line nodes; never uniform) only if profiling demands
it; cross-check against CoolProp.

\textbf{C. First transport closures.} McAdams μ\_m in
Darcy--Weisbach/Churchill (baseline), MSH as sensitivity; regime-logic
wall HTC --- Dittus--Boelter/Gnielinski + natural-convection floor; \citet{shah2017} or Chen superposition + pool-boiling floor for boiling walls; Shah
condensation for cold walls; smooth blends, no hard switches; no
effective k\_m anywhere except an acknowledged bracketing exercise;
±50\% HTC uncertainty band declared up front.

\textbf{D. HLLC retained? Yes.} Max-of-neighbors sound speeds at
dome-straddling faces, positivity-checked intermediate states with local
HLL/Rusanov fallback, post-stage flash with bracketed phase detection,
accepted first-order behavior at the kink. FORCE-style extra dissipation
is the documented fallback if robustness problems emerge, at modest
accuracy cost for a plateau-dominated problem.

\textbf{E. Key HEM limitations (ranked).} (1) Choke discharge
underprediction 0--26\% --- corrected at the boundary, Level 1. (2)
No-slip void/gravity-head error at low rates --- structural, Level 3.
(3) Plateau/undershoot bias from delayed phase change --- Level 2. (4)
Equilibrium condensation onset (supercooled vapor unrepresented) ---
uncharacterized, experiment-facing. (5) No solid phase --- confine to
the choke submodel. (6) Near-critical flash conditioning ---
implementation care, not physics.

\textbf{F. Validation datasets.} SINTEF 2020 Test 3 + N₂ reference + dense tests \citep{munkejord2020,sintef2020data}; \citet{botros2016}; \citet{hammer2022} orifice data; COOLTRANS vents; Guo Test 1; CO₂--N₂/He \citep{munkejord2021,sintef2021data} optional; vessel blowdown for long-time energy balance once \citet{gebbeken1996} or the INERIS sphere is verified/retrieved.

\textbf{G. Quantitative triggers to move beyond HEM.} → Level 1 choke
correction: immediately (evidence already exists). → Level 2 HRM: when
choke-inlet states are two-phase \emph{and} measured discharge or
near-wellhead pressure undershoot disagrees with HEM beyond the C\_d
uncertainty (\textasciitilde±10\%); or when predicted vapor supercooling
(from a passive CNT/supersaturation diagnostic run inside HEM) exceeds a
few K over resolvable durations. → Level 3 drift-flux: when the local
mixture velocity falls below \textasciitilde3× the Shi/Kutateladze drift
velocity (≈1--2 m/s in-bore) over a persistent two-phase interval ---
equivalently when the choked flux drops to the point that G·A\_bore
implies superficial gas velocities under \textasciitilde1 m/s --- or
when measured DTS phase-front positions/hydrostatic gradients contradict
no-slip reconstruction. → Level 4 two-fluid: only on local void/slip
measurements that drift-flux cannot fit.

\textbf{H. Publication opportunities (strongest first).} (1)
\textbf{First vertical, wellbore-scale CO₂ blowdown dataset} ---
distributed P/T + DTS on a 610 m well; Lindeberg called for exactly this
in 2011 and nobody has done it. (2) \textbf{First published CO₂-well
saturation-thermometer DTS trace} --- flashing-onset kinks, saturation
plateaus, phase-front migration; a standalone observational paper. (3)
\textbf{Condensation-side relaxation time for CO₂} --- the fitted θ
literature is all flashing-side; a gas-phase well blowdown with DTS
supercooling signatures could bound it for the first time (complementing \citealp{hammer2025}). (4) \textbf{Transport-closure performance at 170
mm} --- every friction/HTC correlation is extrapolated ≥6× in diameter;
even bulk-calorimetric validation at well scale is novel. (5) The model
paper itself: HEM + conjugate radial network + non-equilibrium choke,
validated on the above --- the composite is unpublished territory.

\textbf{Roadmap.} \emph{Implement now:} two-phase flash (reduced UV,
saturation splines) + dome-aware HLLC guards; McAdams/Churchill
friction; regime-logic HTC with floors; choke saturation-bracketing plus
HF/D-HEM option with C\_c ≈ 0.75 and documented C\_d policy;
verification suite (conservation, stationarity, single-phase limit,
Kumar/Log shock tubes); validation vs SINTEF Test 3 and Botros; emit
T\_sat, supercool margin, x, α, phase-front depths, synthetic DTS.
\emph{Sensitivity study:} MSH vs homogeneous friction;
Cicchitti/Beattie--Whalley μ\_m; Shah vs Chen+floors HTC ±50\%; C\_d ∈
{[}0.74, 0.84{]} × \{HEM, HF, D-HEM\}; θ ∈ {[}0.05, 3{]} ms via a
shadow-HRM run; grid/CFL. \emph{Defer until evidence:} HRM as a
production stage (pending condensation-θ evidence); drift-flux slip
(pending low-rate trigger or DTS contradiction); two-fluid (pending
local data); solid-CO₂/triple-point flash (choke submodel only, pending
observed near-triple-point operation); full (ρ,e) dome tables (pending
profiling).

\hypertarget{conclusion}{%
\subsection{Conclusion}\label{conclusion}}

The critical test the hypothesis was set leaves HEM standing for a
reason worth internalizing: the physics that HEM gets wrong is
concentrated where residence times are short (the choke) and where
velocities are low (the late-blowdown tail), and both failure modes can
be handled without touching the hyperbolic core --- one as a
boundary-condition correction with direct 12.7 mm experimental
precedent, the other as an algebraic drift-flux closure with a
quantitative activation trigger. The deeper strategic finding is that
the validation landscape is inverted relative to the usual situation:
the model is more constrained than the application. Horizontal pipes,
shock tubes and vessels over-determine the EOS, wave dynamics and phase
kinetics, while the application's defining features --- verticality, 170
mm bore, conjugate wall heating, choked finite-inventory drainage, and
the DTS observable --- have \textbf{zero} published data among them.
That makes the planned experiment not a validation exercise for the
model but the field's first dataset in four distinct gaps, and it argues
for designing the model's outputs (supercool margins, phase-front
depths, synthetic DTS) around what the experiment can uniquely measure,
rather than around what the pipe-decompression literature already knows.
